\section{Results}
\label{sec:results}

Tables report throughput (req/s, higher better) and tail latency (p99~ms, lower
better) per access layer and engine. The two consequential patterns are in the
main text (deep fetch, Table~\ref{tab:deep_fetch}; insert,
Table~\ref{tab:write}); the point read, keyset range scan, and aggregation are in
Supplement Tables~S12, S13, and~S15, all from the run of
Section~\ref{sec:methodology}. Each native driver has an untuned policy-selected row
(\texttt{pg}, \texttt{mysql2}) and tuned (\texttt{pg-tuned},
\texttt{mysql2-tuned}), the tuned rows marking a hand-tuned reference point, not one
attained without code changes. We compare \emph{relative}
differences within an engine. Two estimands recur: the \emph{performance of the selected
configured treatment} (RQ1--RQ3), which is not library-only overhead, and the
\emph{common-SQL raw-path sensitivity contrast} (Supplement Table~S14), the residual spread
when every layer runs common SQL through its raw facility (Section~\ref{sec:threats}).

% One combined table per access pattern (throughput + p99), per layer and engine.
% \tabledir is set by each main file so this shared section resolves without TEXINPUTS.
\input{\tabledir deep_fetch}
\input{\tabledir write}

\subsection{RQ1: performance difference relative to the native baseline}
The native implementation leads nine of the ten engine--pattern comparisons:
\texttt{pg} or \texttt{pg-tuned} is fastest on all five PostgreSQL patterns, and
\texttt{mysql2} on four of five MySQL patterns. The exception is the MySQL insert,
where the four fastest treatments span 1{,}766--1{,}816~req/s with broadly overlapping
intervals and the native driver ranks fourth (Table~\ref{tab:write}); that ordering is
practically unresolved rather than a defeat, and we do not read a leader from it. The current-generation ORMs sit mid-to-low on the read patterns; Prisma is near the bottom of point reads and aggregation on both engines (Table~\ref{tab:deep_fetch}; Supplement Tables~S12 and~S15), retiring an earlier version's headline of Prisma tying the native driver at a large CPU premium (Section~\ref{sec:discussion}).

\texttt{pg-tuned} and \texttt{mysql2-tuned} track their untuned counterparts within about
$2\%$ on the deep fetch (\texttt{pg-tuned} 3{,}700 vs.\ \texttt{pg} 3{,}638;
\texttt{mysql2-tuned} 2{,}398 vs.\ \texttt{mysql2} 2{,}416). The one material gain is
PostgreSQL aggregation, where named prepared statements add $11\%$ (7{,}583 vs.\ 6{,}807).
On MySQL the binary-protocol baseline is \emph{slower} than \texttt{mysql2} on every pattern
(by $0.4$--$4.7\%$), so \texttt{execute()} buys nothing at this operating point; both tuned
baselines are therefore disclosed reference points, not optimization ceilings.

On the \textbf{point read} the native-relative spread is $2.83\times$ (PostgreSQL)
and $2.14\times$ (MySQL) (Supplement Table~S12). On the \textbf{deep fetch}
(a post, its author, and all comments with their authors) it is the widest
measured --- $7.01\times$ on PostgreSQL (\texttt{pg} 3{,}638, 95\% CI
[3{,}602--3{,}662], down to MikroORM's 519~[512--528]) and $5.12\times$ on MySQL
(\texttt{mysql2} 2{,}416~[2{,}413--2{,}425] down to MikroORM's 472~[464--475]). The
native driver leads on both engines, the data-mapper ORMs (especially
MikroORM) trailing furthest (Table~\ref{tab:deep_fetch}); the comparison is paired
(Section~\ref{sec:methodology}), every adjacent MySQL step large relative to its
bootstrap interval (Supplement Table~S36). \textbf{Aggregation} is the least
layer-sensitive pattern: the native-relative spread is only $1.82\times$
(PostgreSQL) and $1.87\times$ (MySQL) (Supplement Table~S15). Every layer forwards
a single query for this pattern, which also avoids the deep fetch's nested-graph
materialization; the design does not isolate either feature as the cause of the smaller spread.

\subsection{Common-SQL raw-path sensitivity contrast}
Server-side capture confirms common statement text and database plans across raw paths on both engines; API level, wire protocol, preparation, round trips, and mapping can still differ (Supplement Table~S42). The deep-fetch native-relative spread is then far smaller --- $1.71\times$ (PostgreSQL) and $2.03\times$ (MySQL) against the selected-configuration $7.01\times$ and $5.12\times$ (Supplement Table~S14). Thus the selected paths differ more than these common-SQL raw paths; the compound contrast assigns no share to any mechanism. A
loading-strategy sensitivity check (Supplement Table~S18) argues against a mere
selected-strategy artifact: forcing the join-selected ORMs (Sequelize, MikroORM) to
select-in makes them slower, and Objection's batched-select-in selected path runs about
$1.6\times$ slower than a single join on MySQL. Objection's result is therefore
sensitive to the documented loading option; this comparison does not isolate a library-only cost.

\subsection{RQ2: backend-stack transfer of the ranking}
Switching the backend changes a bundle that can include the DBMS, supported adapter,
lower-level driver, wire protocol, and defaults. RQ2 therefore concerns
\emph{backend-stack transfer}, not a causal DBMS effect. In the primary campaign,
the seven portable layers have stack-ratio heterogeneity on every pattern (blocked
interaction $F=127.1$--$458.5$, permutation $p<5{\times}10^{-5}$, the $1/(B{+}1)$ resolution floor at $B=20{,}000$; Supplement Table~S28),
but this within-campaign test does not establish durable rank reversals.

A second same-host campaign re-runs 25 repetitions of the seven portable layers for deep
fetch, aggregation, and insert under an independently drawn block order (Supplement
Table~S46). It re-randomizes execution order only: the identifier generator is seeded on
endpoint and replicate, so the campaign replays the same draws and cannot detect
sensitivity to the workload sample.
Deep-fetch ranks reproduce exactly on both stacks (Spearman $\rho=1.00$, maximum
median drift $2.3\%$); aggregation reproduces 5/7 PostgreSQL and 7/7 MySQL positions
($\rho=0.96$ and $1.00$). MySQL insert is less stable: only 3/7 positions reproduce,
$\rho=0.82$, the leader changes from Drizzle to Knex, and maximum between-campaign median drift reaches
$13.6\%$. We therefore promote a cross-stack reversal only when it occurs in both
campaigns, exceeds the predeclared $5\%$ practical margin on both stacks, and paired
bootstrap intervals exclude equality in opposite directions in both campaigns.
Four pairs meet all criteria: on deep fetch, Prisma exceeds Objection and Sequelize
on PostgreSQL but trails both on MySQL; on aggregation, Prisma exceeds Objection on
PostgreSQL but trails it on MySQL; on insert, Prisma exceeds MikroORM on PostgreSQL
but trails it on MySQL. All four therefore involve one treatment, and that treatment is
the one whose lower-level driver differs between the stacks: Prisma runs the MariaDB
adapter on MySQL because it ships no \texttt{mysql2} adapter, while every other layer
uses \texttt{mysql2} (Section~\ref{sec:threats}). The recurrence is consistent with
backend-stack sensitivity as RQ2 defines it, but it cannot be separated from that driver
substitution, and it is not evidence about ORMs as a group. Other changes in close ranks,
including the identity of the MySQL insert leader, remain exploratory.

The magnitude of stack transfer is also treatment-specific. Primary-campaign
PostgreSQL$\div$MySQL ratios range from $1.10$ to $1.84\times$ on deep fetch,
$1.27$ to $1.82\times$ on aggregation, and $1.55$ to $3.06\times$ on insert
(Supplement Table~S28). These ratios and the recurring reversals show that the
observed ordering cannot be assumed to transfer unchanged between the two tested
stacks. They do not allocate the difference to the DBMS, adapter, driver, protocol,
or defaults.

The primary MySQL insert medians illustrate why fine ordering is not a headline:
Drizzle, Objection, Knex, and the native driver lie between 1{,}766 and
1{,}816~req/s, while several replicate distributions are visibly multimodal or
heavy-tailed and reach $13.6\%$ within-campaign CV --- numerically equal to the between-campaign drift above, but a different quantity (Supplement Figure~S1). Prisma is
the slowest policy-selected treatment at 894~req/s, giving a $1.98\times$
native-relative spread, but the broad interval on that spread is
[1.65--2.09]. Under relaxed durability the native MySQL and Knex paths gain
$2.42\times$ and $2.07\times$, whereas gains differ by treatment; concurrent
wait-event sums are consistent with a shared commit-path contribution but are not
a causal time fraction (Supplement Tables~S3 and S19). PostgreSQL insert spans
$3.18\times$ (6{,}018 to 1{,}894~req/s) and changes by at most $14\%$ in the
relaxed-durability sensitivity. These compound interventions do not decompose
backend-stack cost.

\subsection{RQ3: which pattern spreads widest}
The primary native-relative spreads on PostgreSQL are deep fetch
($7.01\times$), keyset range scan ($4.22\times$), insert ($3.18\times$), point
read ($2.83\times$), and aggregation ($1.82\times$). On MySQL they are deep fetch
($5.12\times$), range scan ($3.88\times$), point read ($2.14\times$), insert
($1.98\times$), and aggregation ($1.87\times$). Thus the tested deep fetch has the
largest spread and aggregation the smallest on both stacks. A full concurrency
sweep places the median knee at 8--32 connections and the median utilization at 50
connections at 98--100\% for every pattern--stack combination; the lowest
individual values are 87--88\% for two PostgreSQL patterns (Supplement Table~S35).
The fixed point is therefore a high-utilization comparison throughout, although it
is not an identical utilization fraction for every layer. An exploratory fan-out
sweep from 0 to 500 comments finds a roughly fixed multiplier rather than a spread
that grows with breadth; it does not vary graph depth, schema, or query mix.

\subsection{Tail latency}
At equal closed-loop demand, deep-fetch p99 spans 18--108~ms on PostgreSQL and
28--121~ms among the policy-selected layers on MySQL (Table~\ref{tab:deep_fetch}). The run is the analysis unit; a
12~s slow-layer run supplies only about 60 observations to its top percentile, so
we use the median of 25 run-level p99 values and show their raw spread. Ten
60-second repetitions preserve the practical ordering, with 12-versus-60-second rank correlation $1.00$ on both engines; the largest median shift is $5.6\%$ on PostgreSQL and $3.8\%$ on MySQL (Supplement Table~S23). This is
endpoint-specific sensitivity evidence, not a general sufficiency claim for a
12-second p99 window.

At the stable nominal 50\% capacity target, coordinated-omission-corrected p99
contracts to 2.3--4.9~ms on PostgreSQL and 1.9--5.5~ms on MySQL; most 70\% cells
remain similarly small (Supplement Tables~S21--S22 and S43). Re-expressing the
targets against the full-sweep maximum places nominal 50\% at 46.6--51.0\% and
nominal 70\% at 65.2--71.5\%; repeated sweep curves widen those ranges only to
46.2--52.2\% and 64.6--73.0\% (S45). The 85\% and 95\% results are unstable
near-saturation sensitivity observations and do not support a convergence claim.
The stable 50\% result shows that the large equal-demand p99 spread contains a
material capacity-and-queueing component; it does not identify its share or imply
a deployment-general sub-saturation bound.

The two operating points also \emph{order the layers differently}, which is the
sharpest consequence of separating them. Between the equal-demand and 50\%-utilization
p99 orderings, Spearman $\rho$ is $0.64$ on PostgreSQL and $0.29$ on MySQL
(Supplement Table~S43). The native driver moves furthest: fastest on the tail at equal
demand on both engines, it falls to third of eight on PostgreSQL and seventh of eight on
MySQL once each layer is held at half its own capacity. On MySQL the reversal is
unambiguous --- \texttt{knex} 1.9~ms [1.8--2.2] against \texttt{mysql2} 4.7~ms [4.0--5.6],
non-overlapping over five runs --- while on PostgreSQL the native and leading portable
ranges overlap (\texttt{pg} 2.6~[2.2--3.0], \texttt{knex} 2.3~[2.3--2.7]), so only the
MySQL inversion is claimed. Reporting a tail at one arbitrary load point would have
recorded the native driver as the tail leader on both engines. This is direct evidence
that the operating-point separation stage changes a conclusion rather than only its
caveats, and it is why capacity and tail are reported as distinct quantities.

\subsection{Measurement stability and effect sizes}
Deep-fetch throughput CV reaches $4.5\%$, much smaller than its multi-fold spreads.
The largest primary throughput CV is $13.6\%$ on MySQL insert, whose raw replicates
are shown because a median and interval can hide regime mixtures. All 90 cells
contain exactly 25 run-level samples and zero timed errors, timeouts, or non-2xx
responses. One pre-warm-up startup failure was retried once and recorded; no timed
sample preceded that retry. The fail-closed roster validation requires 90$\times$25
observations before accepting the file.

Predeclared native-driver contrasts are foregrounded; ranking-selected adjacent
comparisons remain descriptive. Paired bootstrap intervals characterize only
within-campaign variation, while the same-host validation campaign is reported
separately. Differences of about 1~ms in p99 and throughput ratios inside the
predeclared $\pm5\%$ practical margin are not promoted as substantive ordering.
